\documentclass[11pt]{article}

\begin{document}
\section{Q-learner}
In the Q-learner model, the rat learns the expected value of each action (visiting each port) directly through experience. At each trial, the rat updates its expectation of visiting that port in the direction of the observed outcome with a Rescorla-Wagner learning rule: $Q_{t+1} \leftarrow Q_t + \alpha \delta$ where $\delta = Q_t-r_t$, and $\alpha$ is a learning rate. There's a bit of fudging referring to this as learning Q-values, which normally refers to (state, action) pairs, whereas we're really only learning a value for each leaf, but as for now we're assuming a 1-to-1 equivalence between actions and resulting states, we can get away with it.
\section{Hidden Markov Model}
In the HMM (Hidden Markov Model), instead of learning an expectation of each port from reward history, the agent infers a hidden global state (or contingency, as we're calling them in the experimental setup). Each state represents a particular set of reward probabilities for all ports. For instance, one particular state might assign probabilities of .2 to the first port, .5 to the second port, .2 to the third port, .8 to the fourth port, .2 to the fifth port, and .2 to the sixth port. The states are \textit{hidden}: that is, the rat cannot directly observe which state it is in, and must infer which global state is most likely to lead to the observed reward history. The model is also \textit{Markovian}: whatever the true state is at any time, it can transition to any other state with a fixed probability, and these transitions are Markovian, i.e. only depending on the current state and not on any prior history.

Because we can't observe the true state, we model a \textit{belief} over the hidden states: a probability of the true state being that state. Technically, we assume a set of states, $\phi$, each of which assigns an \textit{emission probability} to each port, which is just the probability of observing a reward at that port. Our belief is represented as a vector $\alpha$, which assigns a weight to each of the states (and thus $alpha$ sums to 1). Our posterior probabilities, equivalent to the Q-values from the earlier model, are given by a weighted average of the emission probabilities, based on the belief state. On each trial, we observe an outcome at one of the ports: the probability of that outcome, given a state, is just the emission probability of that port, if rewarded, and 1 - the emission probability if unrewarded. This gives us a simple update rule for our belief in each possible state, multiplying our previous belief in that state by the probability of the current observation, and then re-normalizing our beliefs to ensure we have a valid probability distribution. Finally, due to the Markovian transitions between states, we assume a small probability of any state transitioning to any other state, which essentially means diffusing our belief state on each trial across all possible states.
\section{Differences}
The primary difference between the two is whether port estimates are independent. In the Q-learner model, each trial provides information about one and only one port. In contrast, for the HMM, each trial provides information about the global state, and thus the estimates of all 6 ports. This is particularly important at contingency switches. In the Q-learner model, a contingency switch will lead to the rat quickly devaluating the two ports it's alternating between on a single stem, but it won't update estimates for the other 4 ports until it switches to a new stem. In contrast, as soon as the HMM believes that the current state, presumably assigning high reward probability to the current stem, is no longer likely, it also increases the belief in high rewards at the other 4 stems.
\end{document}
